\section{Methods: empirical benchmark}
\label{sec:casestudy}

The study evaluates a fixed benchmark of daily PM$_{10}$ point forecasts.
The unit of analysis is a station--model--horizon cell, and the analysis
is deliberately post-evaluation: it audits how eligibility changes when
dynamic-fidelity requirements are added to an error-based screen. This
design keeps the empirical question separate from any attempt to optimise
a forecasting model.

\subsection{Data, forecast task and validation}

The frozen evidence comprises daily PM$_{10}$ series from 17 Spanish
monitoring stations identified in MITECO/EMEP metadata. The stations span
urban, traffic, industrial, suburban, and rural settings. The recovered
station package does not establish one defensible common calendar endpoint
for the complete set, so we report the verified station coverage and daily
resolution rather than infer a period label. The complete grid contains
17 stations, five empirical model families, and seven forecast horizons,
giving 595 structured rows. These rows share stations and underlying time
series; they are evaluation cells, not independent replicates.

The task is direct multi-step daily point forecasting at horizons
$h=1,\ldots,7$. Evaluation follows five expanding rolling-origin folds in
chronological order. For each origin, information that became available
after the origin was excluded from fitting, feature construction,
preprocessing, threshold estimation, and verification inputs. Missing
verification targets were excluded rather than fabricated. The frozen
aggregate does not recover exact per-fold row counts or one complete
train/test calendar, and those details are therefore not asserted.

\subsection{Forecasting models and reference forecast}

The benchmark contains five empirical model families: direct Histogram
Gradient Boosting (HGB), direct Ridge regression, fixed-order SARIMA,
seasonal naive, and STL+Ridge. HGB and Ridge produce direct forecasts for
each horizon. SARIMA is treated as a recursive statistical family,
seasonal naive as a seasonal carry-forward reference, and STL+Ridge as a
decomposition/regression hybrid. The model evaluation specification table
records the inputs and the
configuration details that can be verified from the frozen provenance.

Exact HGB, Ridge, and STL settings are not recoverable from the frozen
release. Historical SARIMA records also conflict; accordingly, we do not
invent an exact order and describe only the fixed-order family supported by
the evaluation package. This limitation affects reconstruction detail,
not the frozen aggregate diagnostics reported here.

Persistence supplies the operational comparator. It is evaluated on the
same matched forecast cases as each candidate model, so positive skill means
that the candidate reduces RMSE relative to persistence on the relevant
cell. The comparator is a reference for the eligibility audit, not a claim
that persistence is the only meaningful baseline.

\begin{table}[ht]
\centering
\scriptsize
\setlength{\tabcolsep}{2pt}
\textbf{Supplementary Table S2. Model evaluation specification.} Reconstructed from the analysis provenance. ``Not recoverable'' denotes a detail not established by the provenance; it is not inferred from model names.\par\smallskip
\begin{tabular}{@{}p{2.15cm}p{2.25cm}p{2.25cm}p{2.05cm}p{2.25cm}p{2.45cm}@{}}
\toprule
Model & Inputs & Forecast strategy & Fold fitting & Verified configuration & Known limitation \\
\midrule
HGB direct & Direct lag-based predictors; complete feature list not recoverable & One direct forecast per horizon & Fit within rolling-origin training data & Family label and 7-horizon evaluation verified & Exact hyperparameters not recoverable from frozen provenance \\
Ridge direct & Direct lag-based predictors; complete feature list not recoverable & One direct forecast per horizon & Fit within rolling-origin training data & Family label and 7-horizon evaluation verified & Exact regularisation and feature list not recoverable \\
SARIMA & Daily PM$_{10}$ series & Recursive statistical forecast & Fit within rolling-origin training data & Fixed-order SARIMA family verified & Historical configuration records conflict; exact order omitted \\
Seasonal naive & Historical seasonal observations & Seasonal carry-forward baseline & No parameter fitting beyond training history & Family label and 7-horizon evaluation verified & Seasonal period details not recoverable from canonical table \\
STL+Ridge & Decomposition and regression inputs & Hybrid decomposition/regression forecast & Fit within rolling-origin training data & Family label and 7-horizon evaluation verified & Exact STL and Ridge settings not recoverable \\
\bottomrule
\end{tabular}
\end{table}

